\documentclass[11pt]{article}
\usepackage[margin=26mm]{geometry}
\usepackage{amsmath,amssymb,amsthm,booktabs,array}
\usepackage[T1]{fontenc}
\usepackage{lmodern}
\usepackage[colorlinks=true,linkcolor=blue,urlcolor=blue,citecolor=blue]{hyperref}
\usepackage{xurl}
\newtheorem{theorem}{Theorem}
\newtheorem{lemma}[theorem]{Lemma}
\newtheorem{corollary}[theorem]{Corollary}
\newcommand{\Z}{\mathbb Z}
\newcommand{\N}{\mathbb N}
\title{Proof of Markey's five-site growth identity\\and Shannon's region formulas}
\author{Haobo Ma\\\small ChronoAI Pte Ltd, Singapore
\and Wenlin Zhang\\\small National University of Singapore\\\small ChronoAI Pte Ltd, Singapore}
\date{November 16, 2026}
\begin{document}
\maketitle
\begin{abstract}
We prove Markey's conjecture in OEIS A400701, FORMULA, revision 30
(October 20, 2026), that the five-site growth count exceeds Shannon's
circle-vertex count at the same iteration $n>1$ by $4n+7$, using
residue-class points in expanding hexagons.
Counting incidences on the same shapes also proves Shannon's conjectured
region formulas and gives the curved-edge counts. The argument proves
the recorded formulas for three and four initial sites and the
generating functions recorded by Irvine.
\end{abstract}

\section{The counts and their common shape}
Use axial coordinates $(x,y)\in\Z^2$ for the triangular lattice, embedded in
the plane as $(x+y/2,\sqrt3y/2)$. Its six neighbor displacements are
\[
 D=\{(1,0),(0,1),(-1,1),(-1,0),(0,-1),(1,-1)\}.
\]
For integers $2\leq m\leq5$, start with $S_m(0)=\{(j,0):0\leq j<m\}$ and set
\[
 S_m(n+1)=S_m(n)\cup\{p: |(p+D)\cap S_m(n)|\geq2\},\qquad
 a_m(n)=|S_m(n)|.
\]
The update is simultaneous. Markey's five-site sequence is $a_5$;
the three- and four-site versions are $a_3,a_4$ \cite{markey,three,four}.
Write $v(n)$ for Shannon's count: retain two distinct initial points and, at
each step, adjoin intersections of distinct circles centered at retained
points, all of radius the initial separation \cite{shannon}.
The totals begin $2,4,8$ at $n=0,1,2$; OEIS lists $v$ from $n=1$.

The conjecture concerns a linear difference between two quadratic counts.
The useful mechanism is a change of coordinates of determinant $3$, not
subtraction of guessed formulas. The formulas and generating functions were
already recorded in OEIS; in particular, Sean A. Irvine supplied the
generating functions for the three-, four-, and five-site sequences in
October 2026 \cite{markey,three,four}. We give direct shape and counting
proofs of these identities. Barker's 2015 petal generating function
already contains the same quadratic arithmetic after a linear adjustment
\cite{petals}; Corollary~\ref{cor:petal} makes that relation explicit,
separately from the geometric petal interpretation.
The shape argument determines the geometry and yields the region and
edge counts. The superficially
similar added-circle sequence A264788 has a different iteration rule
\cite{circles}; its value at stage $3$ is $4$, whereas the present
construction adds $v(3)-v(2)=6$ circles. These statements concern deterministic finite-seed growth,
not a random initial configuration.

Gravner and Griffeath's shape theory treats box-neighborhood growth;
their explicit recursion in Section~2, equation~(2.3), is for the
eight-neighbor Moore rule with threshold three and a diamond seed
\cite{gravner}. Lehjef studies threshold-three percolation on finite
hexagonal arrays \cite[Sections~3.2 and 4]{lehjef}; her residue-class
formulas give completion times for size-dependent seeds, rather than
occupied-site counts from a fixed collinear seed. Neither source contains
Markey's difference identity or the exact two- through five-site
threshold-two growth formulas proved here. Barker's earlier arithmetic
antecedent remains the generating function described above.

\begin{theorem}\label{thm:main}
For $n\geq2$, put $L=m-1+2n$. The map
\[
 (x,y)\longmapsto(u,v)=(x-y+n,x+2y+n)
\]
is a bijection from $S_m(n)$ to
\begin{equation}\label{eq:polygon}
 P(L,n)=\{(u,v)\in\Z^2:0\leq u,v\leq L,\quad
 n<u+v<2L-n,\quad u\equiv v\pmod3\}.
\end{equation}
Moreover $v(n)=a_2(n)$ for all $n\geq0$. Consequently
\begin{equation}\label{eq:difference}
 a_5(n)=v(n)+4n+7\qquad(n\geq2).
\end{equation}
\end{theorem}

\section{Why the hexagon grows by one layer}
For integers $A,L$ let
\[
 H(A,L)=\{(u,v)\in\Z^2:0\leq u,v\leq L,\quad
 A\leq u+v\leq2L-A\}.
\]
The transformed neighbor moves are
\[
 E=\{(1,1),(-1,2),(2,-1),(-1,-1),(1,-2),(-2,1)\}.
\]
\begin{lemma}\label{lem:boundary}
Suppose $A\geq3$ and $L\geq A+2$. A point $p$ is in
$H(A,L)$ or has at least two distinct $E$-neighbors there if and only if
$p+(1,1)\in H(A+1,L+2)$.
\end{lemma}
\begin{proof}
Write $p=(u,v)$. The expanded inequalities, in the old coordinates, are
\begin{equation}\label{eq:expanded}
 -1\leq u,v\leq L+1,\qquad A-1\leq u+v\leq2L-A+1.
\end{equation}
They hold for every old point. If a point violates, for example, $u\geq-1$,
only the move $(2,-1)$ can take it into $u\geq0$; hence it cannot have two
old neighbors. For $u\leq L+1$, $v\geq-1$, $v\leq L+1$, the corresponding
unique possible correcting moves are $(-2,1),(-1,2),(1,-2)$.
If $u+v<A-1$, only $(1,1)$ can reach $u+v\geq A$; the upper sum bound
uses $(-1,-1)$. This proves necessity of every inequality in
\eqref{eq:expanded}.

For sufficiency, suppose \eqref{eq:expanded} holds but $p\notin H(A,L)$.
At least one line of the following table applies. Select the first
applicable line; the two listed moves send $p$ into $H(A,L)$.
\begin{center}\small
\begin{tabular}{lll}
\toprule
Boundary & Additional condition & Two moves\\
\midrule
$u+v=A-1$ & $u\geq1$ & $(1,1),(-1,2)$\\
 & $u\leq0$ & $(1,1),(2,-1)$\\
$u+v=2L-A+1$ & $u\leq L-1$ & $(-1,-1),(1,-2)$\\
 & $u\geq L$ & $(-1,-1),(-2,1)$\\
$u=-1$ & $v\leq A+1$ & $(2,-1),(1,1)$\\
 & $v>A+1$ & $(2,-1),(1,-2)$\\
$u=L+1$ & $L-v\leq A+1$ & $(-2,1),(-1,-1)$\\
 & $L-v>A+1$ & $(-2,1),(-1,2)$\\
$v=-1$ & $u\leq A+1$ & $(-1,2),(1,1)$\\
 & $u>A+1$ & $(-1,2),(-2,1)$\\
$v=L+1$ & $L-u\leq A+1$ & $(1,-2),(-1,-1)$\\
 & $L-u>A+1$ & $(1,-2),(2,-1)$\\
\bottomrule
\end{tabular}
\end{center}
Here the two sum-boundary cases take precedence; otherwise
$A\leq u+v\leq2L-A$. For example, on $u=-1$ we then have
$A+1\leq v\leq L+1$. The move $(2,-1)$ gives $(1,v-1)$, whose sum is $v$;
the second move is $(1,1)$ at $v=A+1$ and $(1,-2)$ at $v\geq A+2$.
Both meet the required coordinate and sum bounds because $A\geq3$ and
$L\geq A+2$. Interchanging coordinates and using
$(u,v)\mapsto(L-u,L-v)$ proves the other three coordinate-boundary cases.
On $u+v=A-1$, the choice $u\geq1$ gives $v\leq A-2$ and
$u\leq A$ by \eqref{eq:expanded}; adding $(1,1)$ or $(-1,2)$ therefore
stays in the coordinate box and raises the sum to $A+1$ or $A$.
If $u\leq0$, use the displayed reflected choice in the second line.
Central reflection handles the upper sum boundary. All listed moves are
distinct. If none of these boundaries occurs, $p$ was already old.
\end{proof}

We now prove the shape assertion. At time $1$ the occupied rows are
\[
 y=-1:\ 1\leq x\leq m-1,\quad
 y=0:\ 0\leq x\leq m-1,\quad
 y=1:\ 0\leq x\leq m-2.
\]
To calculate time $2$, a point $(x,y)$ has neighbors in these rows at
$(x\pm1,y)$, $(x,y+1)$, $(x-1,y+1)$, $(x,y-1)$ and $(x+1,y-1)$.
Testing the row endpoints gives the following complete list:
\begin{center}
\begin{tabular}{c|ccccc}
$y$ & $-2$ & $-1$ & $0$ & $1$ & $2$\\
\hline
$x$ & $[2,m-1]$ & $[0,m]$ & $[0,m-1]$ & $[-1,m-1]$ & $[0,m-3]$
\end{tabular}
\end{center}
An interval with lower endpoint above its upper endpoint is empty.
There are no other rows, since each move changes $y$ by at most one.
Substitution in \eqref{eq:polygon} with $n=2$ gives exactly these intervals:
indeed $3y=v-u$ and $L=m+3\leq8$, so $|y|\leq2$.
At time $n$, write $A=n+1$, $L=m-1+2n$. The change to time $n+1$ shifts
transformed coordinates by $(1,1)$. Lemma~\ref{lem:boundary} applies,
since $n\geq2$ and $L-A=m+n-2\geq2$.
It preserves the congruence $u\equiv v\pmod3$, as every move in $E$
preserves it. Induction proves the shape assertion. The inverse map is
$y=(v-u)/3$, $x=(2u+v)/3-n$, proving bijectivity.

\section{Circle intersections stay on the lattice}
A similarity sends any two initial points to $(0,0)$ and $(1,0)$ and their
radius to $1$. For an axial displacement $(r,s)$, the squared distance is
$Q(r,s)=r^2+rs+s^2$. Distinct lattice centers can have intersecting unit
circles only if $0<Q(r,s)\leq4$. From
$Q=(r+s/2)^2+3s^2/4$ we have $|s|\leq2$; substitution of the five possible
values of $s$ shows that the centers' displacement is, up to the lattice
rotation $(r,s)\mapsto(-s,r+s)$, one of $(1,0),(1,1),(2,0)$.
For centers $0,d$, the unit-circle intersection points in axial coordinates
are respectively
\[
\begin{array}{c|c}
d&\text{intersections}\\\hline
(1,0)&(0,1),(1,-1)\\
(1,1)&(1,0),(0,1)\\
(2,0)&(1,0).
\end{array}
\]
These lists follow either by subtracting the two circle equations or by
substitution followed by the fact that two distinct circles have at most
two intersections (one in the tangency case).
Thus all new intersections are lattice points and are adjacent to both
centers. Conversely a lattice point adjacent to two occupied centers is
such an intersection. Induction identifies the whole circle construction
with $S_2(n)$ and proves $v(n)=a_2(n)$.

\section{The count and its consequences}
Put
\[
 b_r(L)=\left\lfloor\frac{L+3-r}{3}\right\rfloor\ (r=0,1,2),\qquad
 C(n)=|\{(u,v)\in\Z_{\geq0}^2:u+v\leq n,\ u\equiv v\pmod3\}|.
\]
There are $b_r(L)$ integers of residue $r$ in $[0,L]$. The two excluded
corners of the square in \eqref{eq:polygon} are disjoint because $n<L$,
and the reflection $(u,v)\mapsto(L-u,L-v)$ identifies them. Hence
\begin{equation}\label{eq:count}
 |P(L,n)|=\sum_{r=0}^2 b_r(L)^2-2C(n).
\end{equation}
Replacing $L$ by $L+3$ increases each $b_r$ by $1$; their sum is $L+1$.
The two corner counts are unchanged, so
\[
 |P(L+3,n)|-|P(L,n)|=2(L+1)+3=2L+5.
\]
Take $L=2n+1$, the value for $m=2$; its increment by $3$ is the value
for $m=5$. This proves \eqref{eq:difference} and Theorem~\ref{thm:main}.
The restriction is essential: $a_5(0)-v(0)=3$ and $a_5(1)-v(1)=9$.

For completeness, the numbers of points on the slices $u+v=3k,3k+1,3k+2$
with $u\equiv v\pmod3$ are $k+1,k,k+1$, respectively: solve
$2u\equiv u+v\pmod3$ with $0\leq u\leq u+v$.
Summing these slices gives
\[
\begin{array}{c|ccc}
n&3k&3k+1&3k+2\\\hline
2C(n)&3k^2+3k+2&3k^2+5k+2&3k^2+7k+4.
\end{array}
\]
Substitute these expressions and $L=m-1+2n$ in \eqref{eq:count}.
After collecting terms in each residue class one obtains:
\begin{corollary}\label{cor:formulas}
For $n\geq1$,
\[
v(n)=\begin{cases}
(3n^2+5n)/3&n\equiv0\pmod3,\\
(3n^2+5n+4)/3&n\equiv1\pmod3,\\
(3n^2+5n+2)/3&n\equiv2\pmod3.
\end{cases}
\]
For $n\geq2$,
\[
a_5(n)=\frac{3n^2+17n+25-2((n-1)\bmod3)}3,
\]
with $a_5(0)=5$, $a_5(1)=13$. For $n\geq0$ the sibling formulas are
\[
a_3(n)=\begin{cases}n^2+3n+1&n>0,\ n\equiv0\pmod3,\\
n^2+3n+3&\text{otherwise},\end{cases}
\qquad
a_4(n)=\frac{3n^2+13n+12+2(n\bmod3)}3.
\]
\end{corollary}
\begin{proof}
Equation~\eqref{eq:count} proves the formulas for $n\geq2$; the displayed
initial rows give $a_m(0)=m$, $a_m(1)=3m-2$, supplying all remaining cases.
\end{proof}
Write $\Delta^2a(n)=a(n+2)-2a(n+1)+a(n)$.
\begin{corollary}\label{cor:second}
The historical period-three patterns in A400699, A400700 and A400701
are given exactly by
\[
 \Delta^2a_3(n)=\begin{cases}6&n\equiv2\pmod3,\\0&\text{otherwise},\end{cases}
 \quad(n\geq1),\qquad
 \Delta^2a_4(n)=\begin{cases}2&n\equiv0\pmod3,\\0&n\equiv1\pmod3,\\4&n\equiv2\pmod3,\end{cases}
 \quad(n\geq0),
\]
and $\Delta^2a_5(n)=2(n\bmod3)$ for $n\geq2$.
\end{corollary}
\begin{proof}
Substitute the three residue classes in Corollary~\ref{cor:formulas}.
The excluded initial values are
\[
 \Delta^2a_3(0)=2,\qquad \Delta^2a_5(0)=2,\qquad
 \Delta^2a_5(1)=0.\qedhere
\]
\end{proof}
The patterns are recorded in the October 2026 entry histories
\cite{markey,three,four}.

The same formulas prove the recorded generating functions. As identities
of formal power series, with $F_m(z)=\sum_{n\geq0}a_m(n)z^n$,
\[
\begin{aligned}
F_3(z)&=\frac{3+z+2z^2-3z^3+5z^4-2z^5}{(1-z^3)(1-z)^2},\\
F_4(z)&=\frac{2(2+z+z^2-2z^3+z^4)}{(1-z^3)(1-z)^2},\\
F_5(z)&=\frac{5+3z+2z^2-5z^3+z^4-2z^5+2z^6}{(1-z^3)(1-z)^2}.
\end{aligned}
\]
To verify them, write $n=3k+r$ in Corollary~\ref{cor:formulas} and use
$\sum t^k=(1-t)^{-1}$, $\sum kt^k=t(1-t)^{-2}$ and
$\sum k^2t^k=t(1+t)(1-t)^{-3}$, with $t=z^3$, inserting the exceptional
initial values. Equivalently, multiplication by
$1-2z+z^2-z^3+2z^4-z^5$ gives the displayed numerators.
Thus the recurrence with coefficients $(2,-1,1,-2,1)$ holds from $n=6,5,7$
for $m=3,4,5$, respectively; the starting indices matter.

\section{Regions, edges, and added circles}
At iteration $n\geq1$, the drawn circles have centers $S_2(n-1)$ and
their retained vertices are $S_2(n)$. A retained center lies on the circle
of an adjacent retained center. Thus vertices of degree two, including the
initial centers, are included. Let $e(n)$ count the curved edges between
consecutive retained vertices on these circles, and let $r(n)$ count
bounded regions. These are A374339 and A374337 \cite{edges,regions}.

\begin{corollary}\label{cor:regions}
We have $e(1)=6$, $r(1)=3$. For $n\geq2$,
\[
e(n)=\begin{cases}6n^2-4n+2&n\equiv2\pmod3,\\
6n^2-4n-2&n\equiv0,1\pmod3,\end{cases}
\]
and
\[
r(n)=\begin{cases}
(15n^2-17n-3)/3&n\equiv0\pmod3,\\
(15n^2-17n-7)/3&n\equiv1\pmod3,\\
(15n^2-17n+7)/3&n\equiv2\pmod3.
\end{cases}
\]
These prove Shannon's three conjectured region formulas, whose
absolute value in the class $n\equiv1\pmod3$ supplies $r(1)=3$.
\end{corollary}
\begin{proof}
Each circle has at least the two intersections with an adjacent circle,
so its edge count is its number of retained vertices. By Section~3,
\[
 e(n)=\sum_{c\in S_2(n-1)}|\{d\in D:c+d\in S_2(n)\}|.
\]
For $n\geq3$, transform the old set to $P(2n-1,n-1)$. Its coordinate
bounds are $0\leq u,v\leq2n-1$ and
$n\leq u+v\leq3n-2$. In the new transformed coordinates a neighbor
has coordinates $(u,v)+(1,1)+d'$, with $d'\in E$.
For the six respective moves, the only excluded old points lie on
\[
 u+v=3n-2,\quad v=2n-1,\quad u=2n-1,\quad
 u+v=n,\quad v=0,\quad u=0.
\]
For example, $d'=(1,1)$ raises the sum by $4$, so only the old upper
sum boundary exceeds the new upper bound $3n+1$; all other bounds hold.
The other moves give the displayed five boundaries by substitution.
Each coordinate boundary contains
$b=\lfloor(2n-1)/3\rfloor-\lfloor(n-1)/3\rfloor$ points.
Each sum boundary contains
$c=\lfloor(n-\rho)/3\rfloor+1$ points, where
$\rho\in\{0,1,2\}$ is the residue of $2n$ modulo $3$.
Indeed the lower sum boundary requires $u\equiv2n\pmod3$ and
$0\leq u\leq n$; reflection gives the upper one.
Boundary intersections are counted for each missing direction, as required.
Consequently
\[
 e(n)=6v(n-1)-4b-2c.
\]
Substituting Corollary~\ref{cor:formulas} gives the edge formulas.
At $n=1$ the two circles have three retained vertices each, giving $6$
incidences. At $n=2$ the circles centered at $(0,0),(1,0)$ have five
retained vertices each, and those centered at $(0,1),(1,-1)$ have four
each, giving $18$.
The unit-neighbor graph of every $S_2(n-1)$ is connected: the seed is
connected and each added point has two old neighbors. Therefore the
circle arrangement is connected. Euler's formula gives
$r(n)=e(n)-v(n)+1$, including at tangencies and vertices of degree two.
Substitution proves the region formulas and the two initial cases.
\end{proof}

\begin{corollary}\label{cor:circles}
Let $c(n)$ count the circles first drawn at stage $n$ of Shannon's
construction, with the two initial circles at $n=0$. Then
\[
 c(0)=2,\qquad c(n)=v(n)-v(n-1)\quad(n\geq1),
\]
where $v(0)=2$, and
\[
 \sum_{n\geq0}c(n)z^n=
 \frac{2+2z^2+4z^4-2z^5}{(1-z)^2(1+z+z^2)}.
\]
\end{corollary}
\begin{proof}
One new circle is drawn at each new center, and centers determine distinct
circles. Thus the first identity follows from the construction. Multiply
$\sum_{n\geq0}v(n)z^n$ by $1-z$ and sum the three formulas in
Corollary~\ref{cor:formulas}; this gives the displayed rational function.
\end{proof}

There is also an exact algebraic link to Pfoertner's October 4, 2026
observation in the history of A400701 \cite{markey}. Define $p(n)$ as the
coefficient of $z^n$ in Barker's recorded A265676 generating function
\cite{petals}:
\[
 P(z)=\frac{z(1+5z+6z^2+11z^3-5z^4)}{(1-z)^3(1+z+z^2)}.
\]
\begin{corollary}\label{cor:petal}
For these coefficients,
$p(n)=3v(n)-6n-5$ for $n\geq1$, and
$p(n+3)=3a_5(n)-2$ for $n\geq2$.
\end{corollary}
\begin{proof}
Writing $V(z)=\sum_{n\geq1}v(n)z^n$, Corollary~\ref{cor:formulas} gives
\[
 V(z)=\frac{4z+2z^3}{(1-z)^2(1-z^3)},\qquad
 P(z)=3V(z)-\frac{6z}{(1-z)^2}-\frac{5z}{1-z}.
\]
Comparison of coefficients proves the first relation.
Multiplication by the common denominator also gives the polynomial identity
\[
 P(z)=z+7z^2+19z^3+43z^4+
 3z^3\bigl(F_5(z)-5-13z\bigr)-\frac{2z^5}{1-z}.
\]
Comparison of coefficients of $z^{n+3}$ proves the second relation.
\end{proof}
This proves the relation for the recorded power series; an independent
geometric derivation of the petal interpretation of $P$ is a separate
question. In particular, no claim about that geometric count is deduced
from a finite prefix.

\section{Coverage and verification}
\begin{center}\small
\begin{tabular}{p{.26\linewidth}p{.64\linewidth}}
\toprule
Source item & Mathematical coverage\\
\midrule
A400701, FORMULA & Difference conjecture and explicit quasipolynomial proved;
Irvine's generating function and the historical second-difference pattern proved.\\
A374338, FORMULA & Proved: three conjectured vertex formulas;
the Euler relation follows from connectedness.\\
A400699, A400700 & Explicit formulas and Irvine's generating functions and
recurrences proved, with the initial indices stated above.\\
Other lattice rules & A338097, A400975, A400002, A400977 use exactly-two
or restricted-neighbor rules; A008458 is a coordination sequence.
Their statements are not covered.\\
A374337, FORMULA & Proved: Shannon's three conjectured region formulas,
including the exceptional value at $n=1$.\\
A374339 & Not posed; proved: the edge quasipolynomial in
Corollary~\ref{cor:regions}; its Euler identity is also established.\\
A264788 & Not covered. Identification with the added circles in
Corollary~\ref{cor:circles} is not asserted in the source;
it fails at stage $3$ ($4$ versus $6$).\\
A265676; A400701 history & Prior formula cited (Barker, 2015);
Pfoertner's relation follows for its coefficients.
The geometric petal count is not covered.\\
Other circle entries & Not covered: A359569 varies the radius;
A371373, A371254 use regular polygons.\\
\bottomrule
\end{tabular}
\end{center}

The proof above covers all four seed sizes. The pinned Lean declaration
\texttt{MarkeyA400701.result} in
\href{https://github.com/the-omega-institute/trureturing/blob/2f8421f6b90d37b4815b887f6c60987553831b20/D5/S3/StatisticalMechanics/Percolation/MarkeyFiveSiteGrowthDifference.lean}{the pinned difference file}
states \eqref{eq:difference} for every pair of distinct Euclidean initial
points and every $n>1$. Its finite-window count uses the same simultaneous
two-neighbor rule on an induced, nonperiodic lattice window.
\texttt{ShannonA374338.result} in
\href{https://github.com/the-omega-institute/trureturing/blob/2f8421f6b90d37b4815b887f6c60987553831b20/D5/S3/Geometry/ShannonFixedRadiusCircleVertices.lean}{the pinned circle-count file}
states Shannon's three formulas. The declaration
\texttt{AxialGrowth.exact\_shape} in
\href{https://github.com/the-omega-institute/trureturing/blob/2f8421f6b90d37b4815b887f6c60987553831b20/D5/S3/StatisticalMechanics/Percolation/TriangularTwoNeighborGrowth.lean}{the pinned growth file}
covers seed sizes $2$ and $5$, for $n\geq2$.
The extensions to seed sizes $3$ and $4$, the region and edge counts,
the added-circle formula and the power-series identities are written
proofs here. These statements are not conclusions of the three cited
Lean declarations. The accompanying \texttt{verify.py} compares exact
lattice evolution, shapes, formulas, incidences and power-series
coefficients on specified finite ranges, and tests the boundary lemma.
These finite comparisons supplement the all-iteration proofs.

\begin{thebibliography}{9}
\bibitem{markey} A. Markey, \emph{OEIS A400701}, NAME and FORMULA,
revision 30 (October 20, 2026); generating function by S. A. Irvine,
October 19, 2026.
\url{https://oeis.org/A400701}; revision history:
\url{https://oeis.org/history?seq=A400701}.
\bibitem{shannon} S. R. Shannon, \emph{OEIS A374338}, NAME, COMMENTS and
FORMULA, revision 21 (November 10, 2026), submitted July 5, 2024.
\url{https://oeis.org/A374338}.
\bibitem{regions} S. R. Shannon, \emph{OEIS A374337}, NAME and FORMULA,
revision 22 (November 10, 2026), submitted July 5, 2024.
\url{https://oeis.org/A374337}.
\bibitem{edges} S. R. Shannon, \emph{OEIS A374339}, NAME and FORMULA,
revision 17 (November 10, 2026), submitted July 5, 2024.
\url{https://oeis.org/A374339}.
\bibitem{circles} K. Ngaokrajang, \emph{OEIS A264788}, NAME, COMMENTS
and FORMULA, revision 20 (February 16, 2025); generating function by
C. Barker, December 10, 2015. \url{https://oeis.org/A264788}.
\bibitem{petals} K. Ngaokrajang, \emph{OEIS A265676}, NAME and FORMULA,
revision 19 (November 5, 2026); generating function by C. Barker,
December 13, 2015. \url{https://oeis.org/A265676}.
\bibitem{three} A. Markey, \emph{OEIS A400699}, NAME and FORMULA,
revision 38 (November 3, 2026); generating function and recurrence by
S. A. Irvine, October 9, 2026. \url{https://oeis.org/A400699}.
\bibitem{four} A. Markey, \emph{OEIS A400700}, NAME and FORMULA,
revision 35 (November 3, 2026); generating function and recurrence by
S. A. Irvine, October 9, 2026. \url{https://oeis.org/A400700}.
\bibitem{gravner} J. Gravner and D. Griffeath,
\emph{Cellular Automaton Growth on $\Z^2$: Theorems, Examples, and Problems},
definitive author version, April 1998, Sections~1--2, especially
Theorem~1 and equation~(2.3).
\url{https://math.ucdavis.edu/~gravner/papers/cagrowth.pdf}.
\bibitem{lehjef} H. Lehjef,
\emph{Investigating Bootstrap Percolation Through Smart Brute-Forcing},
capstone report, School of Science and Engineering, Al Akhawayn University,
Spring 2023, Sections~3.2 and 4, printed pp.~13--14 and 24--36.
\href{https://cdn.aui.ma/sse-capstone-repository/pdf/spring-2023/INVESTIGATING\%20BOOTSTRAP\%20PERCOLATION\%20THROUGH\%20SMART\%20BRUTE-FORCING.pdf}{Institutional full text}.
\end{thebibliography}
\noindent\small\copyright\ 2026 Haobo Ma and Wenlin Zhang.
Licensed under Creative Commons Attribution 4.0 International (CC BY 4.0),
\url{https://creativecommons.org/licenses/by/4.0/}.
\end{document}
